\documentclass[10pt,a4paper]{amsart}
\usepackage[margin=19mm]{geometry}
\usepackage{amsmath,amssymb,mathtools}
\usepackage[scr=rsfs,cal=euler]{mathalfa}
\usepackage{xfrac}
\usepackage[unicode,hidelinks,bookmarks=false]{hyperref}
\newcommand{\ie}{i.e.\ }
\newcommand{\cf}{cf.\ }
\newcommand{\resp}{respectively\ }
\newcommand{\Subsection}{\subsection}
\providecommand{\nobreakdash}{\nobreak\hbox{-}}
\setlength{\emergencystretch}{2em}
\multlinegap0pt
\usepackage{amssymb}
\usepackage{xcolor}
\newtheorem{lemma}{Lemma}
\newcommand{\dd}{\,\mathrm{d}}
\title{A nonnegativity-preserving finite element method for a class of parabolic SPDEs with multiplicative noise}
\author{Ana Djurdjevac, Claude Le Bris, Endre S\"uli}
\date{Source: arXiv:2502.16854v2 (CC BY 4.0)}
\begin{document}
\maketitle

\noindent\emph{Source and licence.} A nonnegativity-preserving finite element method for a class of parabolic SPDEs with multiplicative noise. Version arXiv:2502.16854v2 (CC BY 4.0), distributed by arXiv under the Creative Commons Attribution 4.0 International licence, http://creativecommons.org/licenses/by/4.0/, as recorded in the arXiv licence metadata for this exact version and shown on the abstract page https://arxiv.org/abs/2502.16854v2. Full licence notice: \texttt{licenses/arxiv-2502.16854-CC-BY-4.0.txt}.\par\smallskip

\noindent\textit{Editorial scope.} Counted unit: the article's lemma L2-h-norm-comparison (source0.tex lines 706--715). The article's complete original proof of it is delivered with it (source0.tex lines 716--754, one \texttt{proof} environment of 39 lines). No mathematics is written, repaired or re-derived: the statement and the proof are the article's own lines, in the article's own notation, and no retained line is substituted or re-flowed.  No further source-local context is needed: every cross-reference of every retained line resolves inside the two delivered spans.  Nothing was omitted from any retained span, so the package carries no omission record.  No retained line contains a citation, so the wrapper carries no bibliography and no bibliography entry.  No retained line contains a figure environment, an image-inclusion command or drawing code, so no image is delivered and the package declares no asset dependency.  Editorial formatting only, in addition: an AMS \texttt{amsart} preamble that restates the article's own declarations and macros used by the retained lines, namely the environments \texttt{lemma} on separate counters, together with the standard packages those lines need.  The retained environments print in this document as Lemma 1 in order of appearance, whereas the article prints the same items with the numbering of its own full document.  The complete native source of the article as distributed in the arXiv e-print for this exact version is bundled unchanged as \texttt{sources/173757/source0.tex}.
\begin{lemma}\label{L2-h-norm-comparison}
There exists a positive constant $C$, independent of $h$, such that, for any function $v_h  \in V_h$ we have
     \begin{equation}\label{eq:17}
                \|v_h\|   \leq \|v_h\|_h \leq C \|v_h\|.  
    \end{equation}
In addition,  for any $v \in C_0(\overline{D})$, 
  \begin{equation}\label{eq:17bis}
  \|\pi_h v\| \leq \|v\|_h.
  \end{equation}
\end{lemma}

\begin{proof}
%
We begin by showing~\eqref{eq:17bis}, which will then, as a special case, imply the first inequality in \eqref{eq:17}.
Suppose that $K$ is a (closed) simplex in the triangulation $\mathcal{T}_h$ of $\overline{D}$, with vertices $x^K_j$, $j=0,\ldots,d$, and $v \in C_0(\overline{D})$. Then, for all~$x \in K$,
%
\begin{align*}
|\pi_hv(x)|^2 &=   \left| \sum_{j=0}^d v(x^K_j) \sqrt{\phi^K_j(x)} \sqrt{\phi^K_j(x)} \right|^2 \\
&\leq \left( \sum_{j=0}^d [v(x^K_j)]^2 \phi^K_j(x)\right) \left( \sum_{j=0}^d \phi^K_j(x)\right)  \\
&= \pi_h[v^2(x)].
\end{align*}
%
Hence, we have that 
%
\begin{align*}
\|\pi_h v\|^2   = \sum_{K \in \mathcal{T}_h}\int_K |\pi_h v(x)|^2 \dd x  \leq    \sum_{K \in \mathcal{T}_h}\int_K \pi_h [v^2(x)] \dd x =   \|v\|^2_h. 
\end{align*}
In particular, with $v= v_h \in {V}_h$, we have  
%
\[ \|v_h\|   = \|\pi_h v_h\|   \leq \|v_h\|_h, \]
%
which is the first inequality in \eqref{eq:17}. 

In order to prove the second inequality in \eqref{eq:17}, we first remark that 
$\|\pi_h [v_h^2]\|_{L^\infty(K)} \leq \|v_h\|^2_{L^\infty(K)}$. Indeed, 
%
\begin{align*}
\|\pi_h[v_h^2]\|_{L^\infty(K)} &= \max_{x \in K} \sum_{j=0}^d v_h^2(x_j^K) \phi_j^K(x)   \leq \max_{j \in \{0,1,\ldots,d\}} v_h^2(x_j^K)\max_{x \in K} \sum_{j=0}^d \phi_j^K(x)\\ &=  \max_{j \in \{0,1,\ldots,d\}} v_h^2(x_j^K) \times 1 \leq \|v_h\|^2_{L^\infty(K)}. 
\end{align*}
We therefore have 
\begin{align*}
\|v_h\|^2_h = \sum_{K \in \mathcal{T}_h} \int_K \pi_h[v_h^2(x)] \dd x & \leq \sum_{K \in \mathcal{T}_h} |K|\, \|\pi_h[v_h^2]\|_{L^\infty(K)} \leq \sum_{K \in \mathcal{T}_h} |K|\, \|v_h\|^2_{L^\infty(K)}\\
&\leq C\,\sum_{K \in \mathcal{T}_h} |K|\, |K|^{-1} \|v_h\|^2_{L^2(K)} = C\|v_h\|^2,
\end{align*}
where the last inequality holds thanks to the assumed shape-regularity of $\mathcal{T}_h$, which implies the existence of  a positive constant $C$ such that, for each closed  simplex $K \in \mathcal{T}_h$,
%
\[\|v_h\|_{L^\infty(K)} \leq C |K|^{-\frac{1}{2}}\|v_h\|_{L^2(K)}.\]
%
This completes the proof of the lemma. 
\end{proof}

\end{document}
